\section{Introduction and Problem Definition}

\subsection{Project Background}
Every semester the timetable committee of an engineering college publishes the class timetable as a
multi-page PDF: one grid per division for each of F.E., S.E., T.E., B.E.\ and M.Tech., with days as
rows, time slots as columns and, below each grid, a legend that expands abbreviations such as
\texttt{DBMS}, \texttt{KKD} or \texttt{LLC}. A single cell may contain a subject, a teacher code, a
room number and, for practical sessions, four stacked lab sub-groups (\texttt{A1}--\texttt{A4}), each
with its own teacher and laboratory. Electives, combined B.E.\ classes, ``Tentative'' sections and
``W.E.F.'' (with-effect-from) dates are written in free text around the grid.

Students and teachers read this document by hand many times a day: to find the next lecture, to see
whether a lab is in 606 or 607, to locate a teacher, or to look for an empty room for a make-up
lecture. When a revised PDF is circulated, everyone has to compare it with the old one manually.
Clashes, for example one professor assigned to two rooms at the same time, are usually noticed only
when they happen in the corridor. Generic calendar apps do not help, because nobody types 450 classes
into a calendar, and timetable-generation software is used only by the committee that creates the
timetable, not by the people who follow it.

\subsection{Problem Statement}
\begin{quote}
\emph{Design and implement a system that converts a college's published timetable (PDF, spreadsheet,
Word file or photograph) into verified structured data, keeps every version as an auditable archive,
detects scheduling clashes, and lets each user (student, professor, HOD, Principal or administrator)
get correct, role-appropriate answers to everyday timetable questions on a mobile phone or laptop,
including questions typed in plain English, without ever exposing one user's personal data to
another.}
\end{quote}

The difficulties with the existing approach are:
\begin{enumerate}
\item \textbf{Unstructured source.} The timetable exists only as a visually formatted document.
      Merged cells, stacked lab rows, page-specific legends, inconsistent time labels (for example a
      row printed as ``12.15--1.15'' without a.m./p.m.) and scanned pages make automatic reading hard.
\item \textbf{Manual lookup.} Answering ``Where is my next class?'' requires locating the right
      page, row, column and lab sub-group, and expanding the legend codes.
\item \textbf{No conflict visibility.} Teacher, room and batch overlaps are not checked mechanically.
\item \textbf{Version confusion.} Old and new PDFs circulate together; it is unclear which one is
      official.
\item \textbf{Privacy.} Students may also want to keep a personal timetable (tuition, coaching),
      which must never be visible to the college administration.
\end{enumerate}

\subsection{Motivation}
The problem affects every member of the college every working day. A typical department timetable
has more than 450 class entries across ten pages, and each of several hundred students and about 50 faculty
members consults it several times a day. Even saving one minute per lookup saves hundreds of
person-hours per semester. More importantly, an automatically checked timetable removes clashes
\emph{before} the semester starts. The project also gave us the opportunity to apply, in one
realistic system, almost every core subject of the first four semesters: data structures and
algorithms in the parser and query engine, DBMS in the schema and transactions, operating-system
concepts in the background worker and concurrency control, computer networks in the client--server
API, and software engineering throughout.

\subsection{Objectives}
The objectives were defined to be specific and measurable:
\begin{enumerate}[label=\textbf{O\arabic*.}]
\item Extract class entries from the college's real timetable PDF with at least 95\% of entries
      verified automatically, and flag every uncertain value for human review instead of guessing.
\item Support at least five input formats (PDF, scanned PDF/photo via OCR, XLSX/CSV, DOCX) through one
      common pipeline.
\item Provide role-based dashboards (current class, next class, free periods) for five roles:
      Student, Professor, HOD, Principal and Admin, with server-side enforcement of permissions.
\item Answer everyday questions typed in English through a deterministic engine covering at least
      25 question types, with clarification instead of guessing when a question is ambiguous.
\item Detect professor, room and batch clashes automatically and report only genuine clashes.
\item Keep institutional and personal timetables strictly separate, with complete version history.
\item Achieve a median search response time below 300~ms for a single user on modest hardware, with
      no incorrect answers under concurrent use.
\item Run on Android, iPhone and laptop browsers from one code base, and deploy a free online demo.
\item Back the system with automated tests (unit, integration, end-to-end) with 100\% of tests passing
      before every commit.
\end{enumerate}

\subsection{Scope of the Project}
\paragraph{Features included.}
Timetable upload and extraction (PDF, scanned PDF, JPG/PNG, XLSX, XLS, CSV, DOCX, DOC); validation,
review and auditable correction; archive of all versions with one official (primary) timetable;
activation notifications; conflict detection; role dashboards; weekly timetable view; natural-language
search (30 intents); free-room and professor-availability search; personal timetables; academic-year
promotion of students; administrator master data (departments, batches, courses, rooms, users,
aliases); audit log.

\paragraph{Features outside the scope.}
Automatic \emph{generation} or optimisation of a timetable (the committee still creates it);
attendance; live location tracking of teachers (the system reports only the \emph{scheduled} room);
exam timetables; integration with the college ERP; free-form chat with a large language model.

\paragraph{Intended users.}
Students, professors, Heads of Department, the Principal and the timetable administrator.

\paragraph{Operating environment.}
Android and iOS phones (as an installable app or a home-screen web app), laptop browsers (Chrome,
Edge, Safari, Firefox) and a Linux server or cloud container running Python 3.12+ and PostgreSQL 14+.

\paragraph{Expected scale.}
One institute: about 2{,}000--3{,}000 users, 50--100 batches, 500--1{,}000 class entries per
timetable, a few timetable versions per semester and short bursts of concurrent use around lecture
changes.

\subsection{Proposed Solution}
The proposed system (Figure~\ref{fig:arch}) has three parts:
\begin{enumerate}
\item \textbf{An extraction pipeline} that converts any supported file into a common \emph{grid}
      model, recognises sections, time rows, lab sub-groups, legends and breaks, validates every entry
      and assigns it a status: \textsc{Verified}, \textsc{Unverified}, \textsc{Incomplete} or
      \textsc{Rejected}.
\item \textbf{A central back end} (FastAPI + PostgreSQL) that stores timetables as archives, enforces
      role-based access on every request, detects conflicts, and answers questions through a
      rule-based natural-language engine with a fixed registry of 30 intents (Q01--Q30) and 22
      parameterised query templates (QT01--QT22).
\item \textbf{A cross-platform app} (React Native/Expo) with role dashboards, a weekly view, an
      ``Ask'' screen, archives, uploads, review screens and administrator tools.
\end{enumerate}

\subsection{Novelty and Originality}
\begin{itemize}
\item \textbf{End-to-end from the real document.} Most student timetable apps expect the timetable to
      be typed in. This system reads the college's actual PDF, including its specific conventions:
      stacked lab sub-groups, page-wise legends, the 12:15 ``noon'' row, combined B.E.\ electives and
      co-curricular LLC slots with no teacher or room.
\item \textbf{Honest uncertainty.} Instead of silently guessing, every entry carries a verification
      status, the raw source text and its page. Ambiguous a.m./p.m.\ is never corrected silently.
\item \textbf{Safe natural-language queries.} Questions are mapped to a closed registry of intents and
      parameterised templates. The engine never builds SQL from user text, which makes it immune to
      prompt or SQL injection by construction and makes every answer reproducible.
\item \textbf{Strict domain separation.} Institutional and personal data live in separate tables;
      the code cannot even express a query that unions them.
\item \textbf{Clash detection with intent awareness.} Parallel lab sub-groups and elective options are
      recognised as intended overlaps, so only genuine clashes are reported (542 raw overlaps reduced
      to the 2 genuine clashes in the real PDF).
\end{itemize}

\subsection{Feasibility Analysis}
\paragraph{Technical feasibility.}
All components use mature open-source software: PyMuPDF for PDF geometry, Tesseract for OCR,
FastAPI and SQLAlchemy for the API, PostgreSQL for storage and Expo for the app. The team had
working knowledge of Python, SQL and JavaScript from earlier semesters. The riskiest part, parsing
the real PDF, was prototyped first and proved feasible (Chapter~10).

\paragraph{Economic feasibility.}
The software stack is free. Development used personal laptops. The online demo runs on Render's free
tier (free web service and a free PostgreSQL database), so the total cost is
\textbf{Rs.~0}. Production hosting for real students is estimated at about
US\$7--25 per month (a small always-on instance and a managed database).

\paragraph{Operational feasibility.}
Users need only a browser or the Expo/installed app. The Admin's workflow (upload, review,
activate) mirrors the existing manual process, and the system never changes the official timetable
without explicit activation.

\paragraph{Time and resource feasibility.}
The work was planned over one semester in phases (requirements, back-end foundation, extraction,
query engine, app, testing, deployment), which fits the mini-project timeline. Section~7.8 shows the
commit history.

\subsection{Social Relevance and Practical Applications}
\begin{itemize}
\item Students, especially first-years and students who switch divisions or electives, can find the
      right class, lab and room instantly, reducing late arrivals and confusion.
\item Faculty can check their own schedule and free periods, and find free rooms for extra lectures.
\item The HOD and Principal get a clash report and a view of room usage, which supports better use of
      classrooms and laboratories, a scarce resource in Indian colleges.
\item The archive gives a transparent record of which timetable was official on which date.
\item The approach transfers to schools, coaching institutes, hospitals (duty rosters) and any
      organisation that publishes grid-style schedules as PDFs.
\item Accessibility: the plain-English ``Ask'' screen, large tap targets and dark mode help users who
      find dense PDF grids hard to read.
\end{itemize}

\subsection{Success Criteria and Key Performance Indicators}
Table~\ref{tab:kpi} lists the KPIs fixed before testing. Chapter~9 (Section~9.10) compares the measured values with
these targets.

\begin{table}[H]
\centering
\caption{Key performance indicators and targets}
\label{tab:kpi}
\small
\begin{tabularx}{\textwidth}{@{}p{2.8cm}X p{3.0cm}X@{}}
\toprule
\textbf{KPI} & \textbf{Definition} & \textbf{Target} & \textbf{Measurement method}\\
\midrule
Extraction verification rate & Verified class entries / all class entries on the real PDF & $\geq 95\%$ & Seed run on the real PDF; regression test\\
False clash reports & Reported conflicts that are not genuine clashes & 0 & Manual check of every reported conflict against the PDF\\
Response time & Median server time per NL search, single user & $< 300$ ms & Timestamped client measurement, 50 requests per query type\\
Throughput & Searches served per second with 4 API workers & $\geq 30$ req/s & Multi-threaded load test (Section 8.6)\\
Reliability & Correct responses / total requests under concurrency & 100\% & Load test, response status check\\
Functional correctness & Passed / total automated tests & 100\% & pytest, Jest, Playwright\\
Extraction time & Time to extract the 10-page real PDF & $< 10$ s & Timed pipeline run (median of 5)\\
Resource usage & Memory of one API worker; database size for one timetable & $< 300$ MB; $< 50$ MB & \texttt{ps}, \texttt{pg\_database\_size}\\
Coverage of questions & Number of supported question types (intents) & $\geq 25$ & Intent registry and its tests\\
\bottomrule
\end{tabularx}
\end{table}

\newpage
